% !TeX root = ../main.tex
% Abstract

Graphics Processing Units, or GPUs, have emerged as powerful parallel computing platforms, offering significant performance advantages for a wide range of applications beyond traditional graphics rendering. To harness this power, compute-focused GPU frameworks have been developed, providing general-purpose parallel programming models. One aspect of these programming models is their memory consistency specification (MCS), which defines the semantics of concurrent shared-memory operations. The MCS also interacts with other aspects of the programming model to describe language-level security guarantees like memory safety, i.e., preventing access to out-of-bounds or uninitialized memory. Understanding when implementations conform to an MCS, and to what degree the MCS provides a sound and useful abstraction of real machines, is crucial for reasoning about GPU programs and for designing effective conformance tests.

This thesis develops several techniques and large-scale studies for validating GPU memory consistency and memory safety. First, it introduces \textit{MC Mutants}, a mutation testing methodology for GPU MCSs that systematically generates mutants of small concurrent GPU programs and evaluates whether testing environments detect them. MC Mutants enables empirical evaluation of MCS test environments, generates 20 conformance tests and 32 mutated tests for WebGPU, uncovers two WebGPU MCS implementation bugs, and leads to the adoption of MCS tests in the official WebGPU conformance test suite. Next, \textit{GPUHarbor} is presented, a browser-based and Android-based tool for large-scale testing of GPU MCSs. By making tests easy to run across many commodity devices, GPUHarbor enables a study of 106 devices from seven vendors, uncovers two new memory consistency bugs, and provides insights into the behavior of diverse GPUs. The lessons learned while developing GPUHarbor also transfer to several other studies, including an analysis of the architectural components of several NVIDIA GPUs and the discovery of a new class of hardware side-channels across vendors based on observations of certain memory consistency behaviors. Finally, \textit{SafeRace} is described, a collection of security assessments and specification proposals aimed at preserving WebGPU's memory safety in the presence of data races. Memory safety in WebGPU is especially important because users regularly execute code from untrusted sources, where a memory isolation failure can expose data from other tabs, processes, or applications. SafeRace assessments are run across dozens of GPUs and 21 WebGPU compilation stacks, identifying memory safety vulnerabilities in GPUs from several vendors, including an AMD vulnerability that was assigned a CVE. SafeRace also identifies a data-race specification vulnerability in WGSL and proposes a validated path toward stronger memory safety guarantees in WebGPU.
